% !TEX root = ../Projektdokumentation.tex
\section{Technisch-wirtschaftlicher Variantenvergleich}
\label{sec:Variantenvergleich}
Für das Projekt wurden drei Lösungsvarianten betrachtet. Die erste
Variante war die Beibehaltung der manuellen Pflege auf Shell-Ebene. Die zweite
Variante war eine Verwaltungsanwendung auf Basis von Django und Django Admin.
Die dritte Variante war die Entwicklung einer schlanken betriebsspezifischen
Webanwendung mit FastAPI und SQLModel. Die Entscheidung wurde sowohl fachlich
als auch wirtschaftlich bewertet.

\subsection{Technische Bewertung der Varianten}
\label{sec:TechnischeBewertung}
Die manuelle Pflege verursacht zunächst keine Einführungskosten, löst aber
keines der beschriebenen Probleme. Der Prozess bleibt weiterhin von einzelnen
Administratoren abhängig. Außerdem fehlen eine saubere Trennung der
Firmen, eine nachvollziehbare Änderungshistorie und eine einheitliche
Prüfung der Eingaben.

Eine Umsetzung mit Django und Django Admin hätte viele Grundfunktionen bereits
mitgebracht. Dazu gehören Benutzerverwaltung, Authentifizierung, Rechteprüfung,
ORM, Formularverarbeitung und eine generische Administrationsoberfläche. Für
das Projektziel war diese Variante jedoch zu umfangreich. Benötigt wurde keine
vollständige Admin-Plattform, sondern eine schlanke Oberfläche für Firmen,
Benutzer, Mail-Adressen, IP-Adressen, Auditierung und Postfix-Export. Django
hätte dafür zusätzliche Framework-Strukturen und Standardmechanismen
eingeführt, die angepasst oder bewusst ausgeblendet werden müssten.

Die schlanke Eigenentwicklung deckt genau die benötigten Funktionen ab und kann
am vorhandenen Ablauf ausgerichtet werden. Die Anwendung bleibt überschaubar,
weil nur die Funktionen umgesetzt werden, die für die Pflege der Cloud-basierten
Kunden benötigt werden. Der Entwicklungsaufwand passt in den geplanten
Zeitrahmen und vermeidet unnötige Framework-Funktionen.

\begin{table}[htbp]
\centering
\begin{tabular}{|p{4.6cm}|c|c|c|}
\hline
\textbf{Kriterium} & \textbf{Manuell} & \textbf{Django Admin} & \textbf{Eigenentwicklung} \\
\hline
Mandantentrennung & schwach & mittel & hoch \\
\hline
Fachprozessabbildung & schwach & mittel & hoch \\
\hline
Nachvollziehbarkeit von Änderungen & schwach & mittel & hoch \\
\hline
Anpassbarkeit & niedrig & mittel & hoch \\
\hline
Betriebliche Einführbarkeit & hoch & mittel & hoch \\
\hline
Einarbeitungsaufwand & leicht & hoch & mittel \\
\hline
\end{tabular}
\caption{Technischer Vergleich der Lösungsvarianten}
\label{tab:TechnischerVergleich}
\end{table}

\subsection{Wirtschaftliche Bewertung}
\label{sec:WirtschaftlicheBewertung}
Die wirtschaftliche Betrachtung orientiert sich an typischen Pflegevorgängen
im Administrationsalltag. Für die Kalkulation wurden 60 Änderungsvorgänge pro
Monat angesetzt. Vor der Projektdurchführung dauerte ein Vorgang im Mittel
rund 10 Minuten, weil Informationen zusammengesucht, Zuständigkeiten geklärt
und Eingaben manuell nachgehalten werden mussten. Mit der neuen Anwendung sinkt
der Aufwand auf rund 4 Minuten, weil die Daten an einer Stelle gepflegt werden
und weniger Rückfragen entstehen.

\begin{table}[htbp]
\centering
\begin{tabular}{|p{8.0cm}|r|}
\hline
\textbf{Kennzahl} & \textbf{Wert} \\
\hline
Pflegevorgänge pro Monat & 60 \\
\hline
Bearbeitungszeit vorher je Vorgang & 10 min \\
\hline
Bearbeitungszeit nachher je Vorgang & 4 min \\
\hline
Zeitersparnis je Vorgang & 6 min \\
\hline
Zeitersparnis pro Monat & 360 min = 6 h \\
\hline
\end{tabular}
\caption{Grundlage der Wirtschaftlichkeitsbetrachtung}
\label{tab:WirtschaftlichkeitGrundlage}
\end{table}

Für die Projektkosten wurde kein Stundenlohn, sondern ein interner
kalkulatorischer Verrechnungssatz angesetzt. Für die 80 Projektstunden wurden
\eur{65} pro Stunde angesetzt. Zusätzlich wurden sechs Stunden fachliche
Begleitung und Abnahme mit \eur{85} pro Stunde bewertet. Daraus ergeben sich
Projektkosten von \eur{5710}. Eine Detailrechnung ist im \Anhang{app:Wirtschaftlichkeit}
dargestellt. Zusätzliche Lizenzkosten fallen für die eingesetzten
Open-Source-Komponenten nicht an. Separate Hardware- oder Hostingkosten wurden
nicht angesetzt, da die Anwendung auf vorhandenen internen Servern betrieben
wird. Der Zugriff ist dabei nicht auf interne Administratoren beschränkt,
sondern auch für autorisierte Firmenbenutzer vorgesehen.

Bei einer monatlichen Einsparung von 6 Stunden und einem kalkulatorischen
internen Satz von \eur{65} ergibt sich ein wirtschaftlicher Nutzen von
\eur{390} pro Monat. Die Projektkosten amortisieren sich damit nach rund
14{,}6 Monaten. Die Amortisationsdauer ist für eine intern genutzte
Verwaltungsanwendung vertretbar. Zusätzlich zur Zeitersparnis werden
Fehlerquellen reduziert, Änderungen besser nachvollziehbar und die
Firmen klarer getrennt.

\subsection{Nutzwertanalyse und Entscheidung}
\label{sec:Nutzwertanalyse}
Da die Bearbeitungszeit allein nicht für die Entscheidung ausreicht, wurde
zusätzlich eine Nutzwertanalyse durchgeführt. Bewertet wurden die Kriterien
Nachvollziehbarkeit, Sicherheit, fachliche Eignung, Bedienbarkeit und
Wartbarkeit. Die gewichtete Auswertung ist im
\Anhang{app:Variantenvergleich} enthalten.

Die manuelle Pflege erzielt den niedrigsten Nutzwert, weil sie den bestehenden
Zustand im Wesentlichen beibehält. Die Django-Variante verbessert zwar die
Bedienung und bringt viele technische Grundlagen mit, wäre für den abgegrenzten
Projektumfang aber zu umfangreich gewesen. Die Eigenentwicklung erzielt den
höchsten Nutzwert, weil sie zum vorhandenen Ablauf passt und ohne unnötige
Zusatzfunktionen eingeführt werden kann.
Aus diesen Gründen wurde die Eigenentwicklung als wirtschaftlich und fachlich
sinnvollste Variante ausgewählt.

Für die Entscheidung war wichtig, dass nicht nur die reine Bearbeitungszeit
betrachtet wurde. Fehlerhafte Freigaben oder unklare Zuständigkeiten können im
Mail-Umfeld direkte Auswirkungen auf den Betrieb haben. Deshalb wurde der
höhere Anfangsaufwand der Eigenentwicklung akzeptiert.
